% Supplementary C — Journal §III (Twelve-axis framework)
% Per-axis motivation + alternative metrics considered + design rationale.

\documentclass[11pt,a4paper]{article}
% Shared preamble for all supplementary chapter documents.
% Used by each supplementary/conference/suppl_*.tex and
% supplementary/journal/suppl_*.tex via % Shared preamble for all supplementary chapter documents.
% Used by each supplementary/conference/suppl_*.tex and
% supplementary/journal/suppl_*.tex via % Shared preamble for all supplementary chapter documents.
% Used by each supplementary/conference/suppl_*.tex and
% supplementary/journal/suppl_*.tex via \input{../preamble.tex}.

\usepackage[utf8]{inputenc}
\usepackage[T1]{fontenc}
\usepackage[a4paper,margin=2.2cm]{geometry}
\usepackage{amsmath,amssymb,amsfonts,amsthm}
\usepackage{mathtools}
\usepackage{graphicx}
\usepackage{booktabs}
\usepackage{array}
\usepackage{multirow}
\usepackage{longtable}
\usepackage{algorithm}
\usepackage{algpseudocode}
\usepackage{xcolor}
\usepackage{url}
\usepackage{cite}
\usepackage[hidelinks]{hyperref}

\graphicspath{{../../figures/}}

\renewcommand{\thesection}{\Alph{section}.\arabic{section}}
\renewcommand{\thesubsection}{\thesection.\arabic{subsection}}

\theoremstyle{plain}
\newtheorem{theorem}{Theorem}[section]
\newtheorem{lemma}[theorem]{Lemma}
\newtheorem{proposition}[theorem]{Proposition}
\theoremstyle{definition}
\newtheorem{definition}[theorem]{Definition}
\theoremstyle{remark}
\newtheorem{remark}[theorem]{Remark}

\setcounter{secnumdepth}{3}
\setcounter{tocdepth}{2}

\newcommand{\R}{\mathbb{R}}
\newcommand{\E}{\mathbb{E}}
\newcommand{\KL}[2]{\mathrm{KL}\!\left(#1 \,\middle\|\, #2\right)}
\newcommand{\ari}{\mathrm{ARI}}
\newcommand{\nmi}{\mathrm{NMI}}
\newcommand{\softmax}{\mathrm{softmax}}
.

\usepackage[utf8]{inputenc}
\usepackage[T1]{fontenc}
\usepackage[a4paper,margin=2.2cm]{geometry}
\usepackage{amsmath,amssymb,amsfonts,amsthm}
\usepackage{mathtools}
\usepackage{graphicx}
\usepackage{booktabs}
\usepackage{array}
\usepackage{multirow}
\usepackage{longtable}
\usepackage{algorithm}
\usepackage{algpseudocode}
\usepackage{xcolor}
\usepackage{url}
\usepackage{cite}
\usepackage[hidelinks]{hyperref}

\graphicspath{{../../figures/}}

\renewcommand{\thesection}{\Alph{section}.\arabic{section}}
\renewcommand{\thesubsection}{\thesection.\arabic{subsection}}

\theoremstyle{plain}
\newtheorem{theorem}{Theorem}[section]
\newtheorem{lemma}[theorem]{Lemma}
\newtheorem{proposition}[theorem]{Proposition}
\theoremstyle{definition}
\newtheorem{definition}[theorem]{Definition}
\theoremstyle{remark}
\newtheorem{remark}[theorem]{Remark}

\setcounter{secnumdepth}{3}
\setcounter{tocdepth}{2}

\newcommand{\R}{\mathbb{R}}
\newcommand{\E}{\mathbb{E}}
\newcommand{\KL}[2]{\mathrm{KL}\!\left(#1 \,\middle\|\, #2\right)}
\newcommand{\ari}{\mathrm{ARI}}
\newcommand{\nmi}{\mathrm{NMI}}
\newcommand{\softmax}{\mathrm{softmax}}
.

\usepackage[utf8]{inputenc}
\usepackage[T1]{fontenc}
\usepackage[a4paper,margin=2.2cm]{geometry}
\usepackage{amsmath,amssymb,amsfonts,amsthm}
\usepackage{mathtools}
\usepackage{graphicx}
\usepackage{booktabs}
\usepackage{array}
\usepackage{multirow}
\usepackage{longtable}
\usepackage{algorithm}
\usepackage{algpseudocode}
\usepackage{xcolor}
\usepackage{url}
\usepackage{cite}
\usepackage[hidelinks]{hyperref}

\graphicspath{{../../figures/}}

\renewcommand{\thesection}{\Alph{section}.\arabic{section}}
\renewcommand{\thesubsection}{\thesection.\arabic{subsection}}

\theoremstyle{plain}
\newtheorem{theorem}{Theorem}[section]
\newtheorem{lemma}[theorem]{Lemma}
\newtheorem{proposition}[theorem]{Proposition}
\theoremstyle{definition}
\newtheorem{definition}[theorem]{Definition}
\theoremstyle{remark}
\newtheorem{remark}[theorem]{Remark}

\setcounter{secnumdepth}{3}
\setcounter{tocdepth}{2}

\newcommand{\R}{\mathbb{R}}
\newcommand{\E}{\mathbb{E}}
\newcommand{\KL}[2]{\mathrm{KL}\!\left(#1 \,\middle\|\, #2\right)}
\newcommand{\ari}{\mathrm{ARI}}
\newcommand{\nmi}{\mathrm{NMI}}
\newcommand{\softmax}{\mathrm{softmax}}


\title{Journal Supplementary C --- Per-axis motivation and alternatives\\
{\large Companion to: \emph{Beyond Accuracy} (Section III)}}
\author{Felipe Santib\'a\~nez-Leal}
\date{2026}

\begin{document}
\maketitle

\section{Scope}
This supplement documents, per axis F-1 through F-12, the question
the axis answers, the metric chosen, the alternatives considered,
and the design rationale for the final selection. The intent is
that a reader could substitute an alternative metric within an axis
without re-architecting the framework.

\section{F-1: Hierarchical Bayesian classification}

\paragraph{Question.} ``Does the topic representation outperform
non-topic baselines after pooling across scenes and folds?''

\paragraph{Metric.} Posterior mean and HDI94 of $\mu_m$ in the
hierarchical model $\text{score}_{m,s,f} = \mu_m + \delta_s + \rho_f
+ \epsilon$, plus the pairwise $P[\mu_A > \mu_B]$ table.

\paragraph{Alternatives considered.} (a)~Pooled OA average (drops
scene-level uncertainty); (b)~Fixed-effect ANOVA (would require
asymptotic normality on small per-fold samples); (c)~Bayesian
paired-tests per (scene, method) (high false-discovery rate without
correction).

\paragraph{Rationale.} The hierarchical model is the smallest design
that (i)~returns a per-method posterior, (ii)~marginalises over
scene-level offsets, (iii)~yields pairwise probabilities directly
interpretable as a graded ordering. The NUTS sampler with 1000
draws + 1000 tune + 2 chains is the default PyMC configuration; no
divergences are observed.

\section{F-2: Topic-word coherence}

\paragraph{Question.} ``Are the top-$N$ words per topic semantically
coherent under the corpus's own co-occurrence statistics?''

\paragraph{Metric.} $c_v$~\cite{roder2015exploring},
$c_{\text{NPMI}}$, and U-Mass, reported with $N = 15$ and a
sliding window of size 110 (matching Röder's default).

\paragraph{Alternatives considered.} (a)~UCI coherence (less
correlated with human judgement on Röder's benchmarks);
(b)~$c_p$ Pointwise mutual information (drops the sliding window);
(c)~Word2Vec embeddings (requires an external corpus, breaks the
self-contained reproducibility commitment).

\paragraph{Rationale.} $c_v$ is the strongest performer in Röder
2015 on human judgements; reporting NPMI alongside is a robustness
check; U-Mass is the historical reference. The three-way report
also enables the coherence-vs-ARI scatter that surfaces the
LDA-vs-ProdLDA tradeoff.

\section{F-3: Seed stability}

\paragraph{Question.} ``Does refitting under a different random
seed yield the same topics?''

\paragraph{Metric.} Mean $\pm$ std ARI(KMeans-on-$\theta$, label)
over $N = 5$ seeds plus mean $\pm$ std $c_v$ over the same five.

\paragraph{Alternatives considered.} (a)~Topic-by-topic matched
cosine (sensitive to permutation, requires pre-alignment);
(b)~Bootstrap-resampled corpora (mixes seed stability with sample
stability); (c)~$N = 30$ seeds (the companion code repository's
\texttt{deep\_seed\_stability/} block runs $N = 30$; the framework
uses $N = 5$ for parsimony plus an explicit upgrade path).

\paragraph{Rationale.} The ARI + $c_v$ pair is the smallest set
that surfaces both the cluster-quality and topic-interpretability
sides of stability. The five-seed budget is enough to estimate the
std reliably for a unimodal posterior; multi-modal posteriors (which
we suspect on KSC and Botswana) would benefit from larger $N$ and
are flagged for follow-up.

\section{F-4: Capacity sensitivity}

\paragraph{Question.} ``How does the topic representation change
when one allows $K$ to vary?''

\paragraph{Metric.} ARI(canonical, KMeans-on-$\theta_K$, label),
held-out perplexity, topic diversity (fraction of unique top-15
words), and cross-seed matched-cosine, all reported across
$K \in \{4, 6, 8, 10, 12, 16, 20, 24\}$.

\paragraph{Alternatives considered.} (a)~ELBO vs $K$ curve
(diverges in batch sizes and not directly interpretable);
(b)~AIC / BIC at the corpus level (penalty terms become arbitrary
for non-Gaussian likelihoods); (c)~Held-out perplexity alone (the
de-facto default but uninformative on topic interpretability).

\paragraph{Rationale.} The four-metric report covers the three
distinct ways capacity choice can fail: underfitting (perplexity
high, ARI low), overfitting on tokenisation noise (perplexity low,
diversity collapses), and topic splitting/merging without ARI gain
(matched-cosine high but ARI flat).

\section{F-5: Band-mask robustness}

\paragraph{Question.} ``Does the topic basis survive a class of
spectral-window restrictions?''

\paragraph{Metric.} Paired ARI between canonical and masked
dominant-topic maps, restricted to pixels labelled in both, under
four canonical masks (VNIR, SWIR, no-water, top-50 Fisher).

\paragraph{Alternatives considered.} (a)~Phi-cosine after Hungarian
alignment (basis-level, but conflates marginal-prevalence with
identity); (b)~KL between $\theta$ posteriors per pixel
(continuous, but interpretation across pixels is unclear);
(c)~Jaccard of top-words per topic after Hungarian alignment
(vocabulary-level, useful supplementary B axis in the conference
paper).

\paragraph{Rationale.} Paired ARI is the metric most directly
interpretable as ``how much agreement under best-permutation'' on
the same data set; it is the principal metric. The companion
swap-rate is the un-corrected fraction and the Hungarian alignment
is the per-topic correspondence; the three together constitute the
axis output.

\section{F-6: Cross-method clustering agreement}

\paragraph{Question.} ``How does the canonical topic-dominant map
relate to other partitions of the same scene?''

\paragraph{Metric.} All-pairs ARI and NMI matrix between $\{$label,
topic-dominant, Felzenszwalb, SLIC-500, SLIC-2000, patch-15,
patch-7, pixel$\}$.

\paragraph{Alternatives considered.} (a)~Pair-by-pair ARI with a
sole reference (loses the cross-method symmetric structure);
(b)~V-measure heatmap (V-measure equals NMI on these partitions
and gives the same information); (c)~Graph-clustering modularity
between methods (interesting but requires a per-method graph
that we do not produce).

\paragraph{Rationale.} The 8$\times$8 matrix surfaces shared
inductive biases (e.g.\ SLIC-500 and patch-15 are spatially smooth
and agree highly even though they disagree with the ground-truth
label) that a single-row reference cannot show.

\section{F-7: Topic--label coupling}

\paragraph{Question.} ``How concentrated is the support of each
canonical topic on the ground-truth label?''

\paragraph{Metric.} Per-topic
$H(L \mid t) = -\sum_\ell P(\ell \mid t) \log P(\ell \mid t)$ and
the corresponding KL divergence $\KL{P(L \mid t)}{P(L)}$.

\paragraph{Alternatives considered.} (a)~Topic-to-label hard
mapping via Hungarian (a special case of F-8 applied to the label
side); (b)~Mutual information $I(L; T)$ (a corpus-level summary,
useful as a single-number companion); (c)~Topic precision/recall on
the dominant-label hypothesis (per-(topic, label) but multiplies
the artefact size).

\paragraph{Rationale.} Per-topic $H$ and KL admit direct
interpretation as concentration / divergence from the marginal,
and integrate cleanly with the dominant-topic-map raster that the
framework already produces.

\section{F-8: Per-topic identity tracking under masks}

\paragraph{Question.} ``Which masked topic corresponds to which
canonical topic, and how much disagreement survives the
correspondence?''

\paragraph{Metric.} Hungarian permutation $\sigma^* = \arg\max_\sigma
\sum_k C_{k, \sigma(k)}$ on the confusion matrix
$C \in \mathbb{Z}^{K \times K^{(m)}}_{\ge 0}$, plus the swap rate
under alignment $|\mathcal{S}|^{-1} \sum_d \mathbf{1}[z^{(m)}_d \ne
\sigma^*(z_d)]$.

\paragraph{Alternatives considered.} (a)~Top-$N$ greedy assignment
(faster but sub-optimal); (b)~Soft alignment via Sinkhorn (returns
a transport matrix, not a permutation); (c)~Spectral alignment of
the eigenvectors of $C$ (numerical stability concerns at small
$K$).

\paragraph{Rationale.} The exact Hungarian solution is computable
in $O(K^3)$ on $K \le 12$ in milliseconds and is the standard
benchmark in the cluster-comparison literature.

\section{F-9: Preprocessing stability (HIDSAG)}

\paragraph{Question.} ``Does the topic basis survive a change in
the preprocessing pipeline (continuum removal, smoothing,
derivative)?''

\paragraph{Metric.} ARI(dominant-topic recipe A, dominant-topic
recipe B) per pair, summarised as the mean ARI across the
$\binom{R}{2} = 3$ pairs for $R = 3$ recipes.

\paragraph{Alternatives considered.} (a)~Phi-cosine pairwise
(basis-level, parallel to F-5's alternative); (b)~Recipe-conditioned
classification accuracy (drops the topic-level question);
(c)~Procrustes alignment of $\phi$ matrices (returns a continuous
similarity but does not produce a discrete per-pair scalar).

\paragraph{Rationale.} Mean pairwise ARI is the simplest summary
that yields a per-subset scalar comparable across HIDSAG subsets.

\section{F-10: Cross-scene topic transfer}

\paragraph{Question.} ``Does the canonical topic basis on scene $s$
generate a useful topic representation on scene $t$?''

\paragraph{Metric.} ARI(label$_t$, dominant-topic of $\theta^{(t)} =
\theta(\phi^{(s)}, \mathbf{c}^{(t)})$) for the transfer
$\phi^{(s)} \to \theta^{(t)}$.

\paragraph{Alternatives considered.} (a)~Held-out perplexity of
$\phi^{(s)}$ on scene $t$ (loses the cluster-quality side);
(b)~Top-words drift between $\phi^{(s)}$ and a refit
$\phi^{(t)}$ (a basis-level statistic that does not say whether the
transfer is useful downstream).

\paragraph{Rationale.} The ARI report is the most directly
interpretable transfer score and is symmetric in the sense of
producing both $s \to t$ and $t \to s$ numbers.

\section{F-11: Rate--distortion of $\theta$}

\paragraph{Question.} ``How aggressively does the topic
representation compress the spectrum while remaining useful?''

\paragraph{Metric.} Per-pixel MSE($\hat{\mathbf{x}} = \phi^\top
\theta, \mathbf{x}$) versus rate $\log_2 K$.

\paragraph{Caveat.} The reconstruction $\hat{\mathbf{x}}$ uses the
quantised band-frequency tokenisation; the MSE therefore mixes a
tokenisation-loss component into the rate--distortion curve. We
flag this as a limitation in journal §VII and as a candidate for
follow-up axis design.

\section{F-12: External baseline comparison}

\paragraph{Question.} ``How does the canonical topic-routed-soft OA
compare to the literature on the same scene?''

\paragraph{Metric.} OA delta = topic-routed-soft OA $-$ best
published OA per scene, where ``best published'' is taken from a
fixed set of three reference papers per scene.

\paragraph{Caveat.} Train/test splits across papers are not always
comparable, so the delta is reported with the per-paper split
caveat in journal §VII and in the artefact.

\bibliographystyle{IEEEtran}
\bibliography{../../bibliography/refs}

\end{document}
